\newpage
\section{Background Knowledge and Technologies}

This chapter assembles the evidence on which the proposed design rests.

\subsection{Experiential learning and realistic job previews}

Kolb's model of experiential learning\cite{r3} characterises learning as a cycle in which concrete experience is transformed through reflective observation into abstract conceptualisation. Meta-analytic evidence associates realistic job previews with increased self-selection and lower voluntary turnover\cite{r4}.

\subsection{Assessment of competence from behaviour}

Evidence-centered design\cite{r5} decomposes the assessment problem into three linked models, summarised in Table 2.1.

{\small
\begin{longtable}{|>{\raggedright\arraybackslash}p{\dimexpr 0.159\textwidth-2\tabcolsep-\arrayrulewidth\relax}|>{\raggedright\arraybackslash}p{\dimexpr 0.317\textwidth-2\tabcolsep-\arrayrulewidth\relax}|>{\raggedright\arraybackslash}p{\dimexpr 0.519\textwidth-2\tabcolsep-\arrayrulewidth\relax}|}
\caption{Evidence-centered design models and the requirements they impose.} \\
\hline
\textbf{Model} & \textbf{Function} & \textbf{Requirement it imposes} \\ \hline
\endfirsthead
\hline
\textbf{Model} & \textbf{Function} & \textbf{Requirement it imposes} \\ \hline
\endhead
Competency & Defines the constructs claimed to be measured & Constructs must be specified independently of any single task \\ \hline
Evidence & Maps observed behaviour to inferences about the constructs & Each observable must be attributable to a named construct \\ \hline
\end{longtable}
}
